\documentclass[addpoints,12pt,answers]{exam}
\usepackage{babel}
\usepackage[top=3cm,bottom=3cm,left=2.5cm,right=2.5cm]{geometry}
\usepackage[utf8]{inputenc}
\usepackage[T1]{fontenc}
\usepackage{booktabs}
\usepackage{graphicx}
\usepackage{csquotes}
\usepackage{paralist}
\usepackage{wasysym}
\usepackage[math]{iwona}
\usepackage{setspace}
\usepackage{enumitem}
\usepackage{multirow}

\pointpoints{Point}{Points}
\bonuspointpoints{Bonuspunkt}{Bonuspunkte}
\renewcommand{\solutiontitle}{\noindent\textbf{Solution:}\enspace}

\chqword{Frage}   
\chpgword{Seite} 
\chpword{Punkte}   
\chbpword{Bonus Points} 
\chsword{Erreicht}   
\chtword{Gesamt}

\checkboxchar{\Square}
\checkedchar{\CheckedBox}

\pagestyle{headandfoot}
%\runningheadrule
%\firstpageheader{ECO208}{Test 3}{28th Feb, 2022}
%\runningheader{ECO208}{Test 3}{28th Feb, 2022}
%\firstpagefooter{ECO208}{Test 3}{\thepage\,/\,\numpages}
%\runningfooter{ECO208}{Test 3}{\thepage\,/\,\numpages}
\doublespacing

\begin{document}
	\vspace*{3em}
	
	\makebox[\textwidth]{ECO 202  Sample Short Answer Question}
	\makebox[\textwidth]{Name:\enspace\hrulefill}
	
	\vspace*{3em}
	
	\begin{questions}
		
		
		\question   Your colleague asserts that a higher multiplier is good because it increases the efficacy of fiscal and monetary policies in combating recessions. Is this statement correct? Explain why or why not.
		
		
		\textbf{Solution}
		Meanwhile, it is correct that the policy would become more effective in combating the recession. Any contractionary policy would also have larger effect as well. If the government has a tight budget and force to reduce government spending, it would trigger a larger recession with higher multiplier.
		
		\question Explain the difference between flow and stock measures of economic, while identifying specific examples. Include a discussion (again on the basis of specific examples) of how such a distinction might be critical for comparing patterns of macroeconomic well-being and growth for different countries, regions and/or individuals.
		
		\textbf{Solution}
		Flow measures account for economic activity over a specified period, typically representing a rate of change or flow per unit of time. Examples of flow measures include: output, investment, etc
		
		Stock measures represent the accumulation of a particular economic variable at a specific point in time. Examples of stock measures include: Capital, debt, wealth
		
		When comparing the economic performance of countries or regions, flow measures like GDP growth rates provide insights into the economic activity over certain period. However, stock measures such as wealth/capital per capita offer a more comprehensive view of economic well-being, it shows the accumulated assets and resources available to individuals or nations.
		
		
		\question  Determine the scale of return of the followings production function where Y is the output, N is the labour input and K is the capital input
		\begin{parts}
			\part $Y= A(N+K)^\alpha$
			\part $Y = (N^\alpha + K^\alpha)^{\frac{1}{\alpha}}$
			\part $Y = AN^{0.3} K^{0.7}$
		\end{parts}
		
		\textbf{Solution}
		For each of the following we will multiple the input by $t$ times
		\begin{parts}
			\part $Y= A(N+K)^\alpha$
				\[Y= A(tN+tK)^\alpha = A t^\alpha (N+K)^\alpha\]
				It will depend on the $\alpha$.
				
				If $\alpha <1$, the function is decreasing return to scale
				
				If $\alpha =1$, the function is constant return to scale
				
				If $\alpha >1$, the function is increasing return to scale
			\part $Y = (N^\alpha + K^\alpha)^{\frac{1}{\alpha}}$
			\[Y = (N^\alpha + K^\alpha)^{\frac{1}{\alpha}} = \left((tN)^\alpha + (tK)^\alpha\right)^\frac{1}{\alpha} = \left((N^\alpha + K^\alpha) (t)^\alpha\right)^\frac{1}{\alpha} = t (N^\alpha + K^\alpha)^{\frac{1}{\alpha}}\]
				This function is constant return to scale
			\part $Y = AN^{0.3} K^{0.7}$
			
			\[Y = AN^{0.3} K^{0.7} = A(tN)^{0.3} (tK)^{0.7} = t AN^{0.3} K^{0.7}\]
			
			This function is constant return to scale
		\end{parts}
		
		\question  List and explain three different ways that the government can finance their spending.
			
		\textbf{Solution}
		\begin{parts}
			\part 	Increasing tax revenue: the government can finance its spending is by raising tax revenue. This involves imposing higher taxes on individuals, businesses, or certain economic activities to generate additional income for the government. 
			\part	Increase money supply: the government finance its spending is by increasing the money supply. This typically involves the central bank purchasing government securities. The government can then use this newly created money to finance its expenditures.
			\part	Increase the amount of government borrowing: increasing the amount of government borrowing. This entails the government issuing bonds or taking out loans from domestic or international sources to raise funds
		\end{parts}
		
		\question  What is Sticky prices? In the stylized DSGE model, how does it lead to inefficiency in the economy?
		
		\textbf{Solution}
		Sticky prices refer to prices that do not adjust immediately in response to changes in economic conditions. For example, the menu cost, wage bargaining etc all could lead to prices be sticky in the short run. 
		In the stylized DSGE model, sticky prices lead to inefficiency in the economy because the supply may not be equal to the demand as the price cannot be adjust such that the market would be clear
		

%		\newpage 
%		\centering{This page intentionally left blank.}
%		\ % The empty page
%		
		\vfill
		
		 \centering{-- The End -- }
		
		
		
		%		\begin{solution}
			%			
			%		\end{solution}	
%		\clearpage 
%		
%		\question Was wäre, wenn es keine Luft gäbe?
%		\begin{parts} 
%			\part[5] Was würde mit Luftballons geschehen? 
%			\bonuspart[6] Wie könnten Fluggesellschaften damit umgehen?
%		\end{parts}
%		
%		\question [100] Wird es morgen schneien?
%		\begin{checkboxes}
%			\CorrectChoice Ja
%			\choice Nein
%			\choice Vielleicht
%		\end{checkboxes}
%		
%		\question Ein Name der folgenden Reihe passt nicht zu den anderen. Welcher?
%		\begin{oneparchoices}
%			\choice Donald
%			\choice Dagobert
%			\choice Daisy
%			\choice Micky
%			\CorrectChoice Balu
%		\end{oneparchoices}
		
	\end{questions}
	
%	\begin{center}
%		\combinedgradetable[h][questions]
%	\end{center}
	
\end{document}